\documentclass[a4paper,11pt]{article}
\usepackage[utf8]{inputenc}
\usepackage[T1]{fontenc}
\usepackage[french]{babel}
\usepackage[left=2cm,right=2cm,top=2cm,bottom=2cm]{geometry}
\usepackage{xcolor}
\usepackage{hyperref}
\usepackage{fontawesome5}
\usepackage{parskip}

% --- Couleurs Dassault Systèmes ---
\definecolor{primary}{RGB}{0, 32, 91}
\hypersetup{colorlinks=true, urlcolor=primary}

\begin{document}

% --- EXPÉDITEUR ---
\noindent
\begin{minipage}[t]{0.5\textwidth}
    \textbf{Loïc Aron MBASSI EWOLO}\\
    Élève-Ingénieur Bac+5 -- Double diplôme ENSEA\\[3pt]
    \faMapMarker\ Cergy, France\\
    \faPhone\ (+33) 7 59 52 33 98\\
    \faEnvelope\ \href{mailto:loicaron.mbassiewolo@ensea.fr}{loicaron.mbassiewolo@ensea.fr}
\end{minipage}
\hfill
\begin{minipage}[t]{0.4\textwidth}
    \raggedleft
    Cergy, le 20 mars 2026
\end{minipage}

\vspace{1.2cm}

\hfill
\begin{minipage}[t]{0.5\textwidth}
    \raggedleft
    \textbf{Dassault Systèmes}\\
    Équipe Generative Experiences -- ENOVIA\\
    Vélizy-Villacoublay, France
\end{minipage}

\vspace{1cm}

\noindent\textbf{Objet :} Candidature au stage -- Data Scientist | Systèmes Multi-Agents IA \& LLM (F/H)

\vspace{0.8cm}

Madame, Monsieur,

Concevoir un système de 5 agents IA collaboratifs orchestrés via Google ADK, avec LangChain pour la chaîne de raisonnement, RAG pour l'enrichissement contextuel, et Gemini comme moteur LLM : c'est exactement ce que j'ai conçu et déployé dans le cadre d'un projet de recherche sur les recommandations agricoles multi-critères. C'est cette expérience concrète en systèmes multi-agents qui m'amène à vous adresser ma candidature pour ce stage au sein de l'équipe Generative Experiences d'ENOVIA.

Le système que j'ai développé répond à plusieurs des défis que vous décrivez : orchestration de plusieurs agents avec gestion de conflits et cohérence des réponses, pipeline RAG pour l'augmentation contextuelle du prompt, intégration d'APIs REST externes (OpenWeather, données de marché), et déploiement production via FastAPI et Docker. La mise en place de patterns de raisonnement inspirés de ReAct (observation, raisonnement, action) pour guider les décisions des agents m'a permis de comprendre les compromis fondamentaux entre qualité de raisonnement, latence et coût d'inférence. LangGraph, que vous mentionnez, représente pour moi l'évolution naturelle de ce travail sur LangChain, et je l'explore activement.

Mon stage de recherche à l'INRIA Saclay (Équipe TRIBE, Institut Polytechnique de Paris, Avril-Août 2026) complète cette expérience sur le plan du NLP à grande échelle : analyse de politiques de confidentialité par BERT/Transformers, conception de pipelines de traitement automatisé, et modélisation de scores interprétables pour des utilisateurs non-experts. Ce travail de recherche m'a appris la rigueur nécessaire pour évaluer et documenter objectivement les performances d'un système IA, ce qui correspond directement à la mission de test et d'optimisation décrite dans votre offre. Par ailleurs, mon rôle d'Assistant de Recherche au MILA (Institut Québécois d'IA, été 2025) m'a confronté à l'étude de l'interprétabilité et du comportement émergent des LLMs, un sujet central pour garantir la fiabilité d'un assistant virtuel industriel.

Je suis disponible dès septembre 2026, à l'issue de mon stage INRIA. Contribuer au développement de l'assistant virtuel multi-agents d'ENOVIA, au coeur de la plateforme 3DEXPERIENCE qui équipe les plus grands industriels mondiaux, représente pour moi une opportunité unique de donner une envergure concrète à mes travaux sur les LLMs et les agents IA. Je serais ravi d'en discuter avec vous lors d'un entretien.

Je vous prie d'agréer, Madame, Monsieur, l'expression de mes salutations distinguées.

\vspace{0.8cm}

\textbf{Loïc Aron MBASSI EWOLO}\\
\textit{Élève-Ingénieur ENSEA / Polytechnique Yaoundé}

\end{document}
